\chapter{Felhasználói élmény és reszponzív dizájn}

Az előző fejezet a WaccApp frontendjének tervezési folyamatát mutatta be: milyen felhasználói folyamatokból, nézetekre bontási döntésekből, adatáramlási elvekből, validációs és állapotkezelési szempontokból épült fel a rendszer. Ebben a fejezetben már az kerül előtérbe, hogy ezek a tervezési döntések hogyan jelennek meg a felhasználói élményben és a reszponzív működésben.

A WaccApp felületének használhatósága azon múlik, hogy a felhasználó mennyire gyorsan érti meg az adott nézet célját, mennyire egyértelmű számára a következő lépés, és kap-e időben visszajelzést a hibás vagy sikeres műveletekről. A UX-fejezet ezért nem újra a szerkezeti döntéseket tárgyalja, hanem azt vizsgálja, hogyan segítik ezek a döntések a mindennapi használatot: az adatbevitelt, a hibajavítást, a beszélgetések olvasását, a szűrést és a kisebb kijelzőn történő kezelést.

\section{Felületi szervezés és információs hierarchia}

A WaccApp felületi szervezésénél az volt a legfontosabb, hogy a felhasználó az adott oldalon gyorsan felismerje, mi a következő értelmes lépés. Egy kommunikációs rendszerben különösen zavaró, ha a felület egyszerre túl sok lehetőséget mutat, de nem derül ki, melyik művelethez melyik mező tartozik. Ezért minden fő nézetnek saját hangsúlyt kellett kapnia.

Az \path{index.html} oldalon a felhasználó küldési folyamatot indít. Itt a címzett, a küldési mód, a tartalom és az indítási lehetőség kerül előtérbe. A felületnek azt kell támogatnia, hogy a felhasználó először kiválassza, milyen típusú kommunikációt szeretne indítani, majd ehhez igazodva adja meg a szükséges adatokat. Egy szöveges üzenetnél a címzett és az üzenettartalom a fontos, sablonos küldésnél viszont már a sablon és a paraméterek is szerepet kapnak. Ha ezeket a mezőket a felület nem választaná el egymástól egyértelműen, a használat könnyen bizonytalanná válna.

A \path{filter.html} egészen más logikát követ. Itt a felhasználó nem új kommunikációt indít, hanem korábbi adatok között keres. Emiatt ezen az oldalon a szűrési feltételek sorrendje és értelmezhetősége fontosabb, mint egy nagy, látványos műveleti gomb. A név, telefonszám, üzenettípus, dátumintervallum vagy kérdőív szerinti keresés akkor használható jól, ha a felület nem egyszerűen egymás alá helyezi a mezőket, hanem segít abban, hogy a felhasználó értelmes keresési feltételt állítson össze.

A kérdőívszerkesztő oldalaknál a legnagyobb kihívást az jelentette, hogy a felhasználó nem egyetlen adatot ad meg, hanem egymással összefüggő elemeket épít fel. Egy kérdéshez válaszok tartozhatnak, a válaszokhoz pedig következő lépések kapcsolódhatnak. Ilyen helyzetben a felületnek nem elég szép beviteli mezőket adnia. Segítenie kell abban is, hogy a szerkesztés közben a folyamat logikája követhető maradjon. Ha a kérdés, a válasz és a következő lépés vizuálisan túl távol kerül egymástól, a felhasználó könnyebben hibázik.

A \path{conv.html} beszélgetésnézete szintén külön felületi szervezést igényel. Itt nem az adatbevitel az elsődleges, hanem az időrendben megjelenő kommunikáció értelmezése. A bejövő és kimenő üzenetek elkülönítése, az időpontok követhetősége és a médiaelemek megjelenítése mind azt szolgálja, hogy a felhasználó ne nyers adatsorokat lásson, hanem egy beszélgetési folyamatot. Ez a nézet akkor működik jól, ha a felhasználó gyorsan megérti, ki küldött üzenetet, mikor történt a küldés, és milyen tartalom kapcsolódott az adott kommunikációhoz.

\section{Reszponzív működés és képernyőmérethez igazított felület}

A WaccApp reszponzív kialakításánál nem az volt a cél, hogy az asztali nézet egyszerűen kisebb méretben jelenjen meg mobilon. Egy kisebb kijelzőn másképp használja az ember a rendszert. Kevesebb információ fér el egyszerre, az érintéses vezérlés miatt nagyobb szerepet kapnak a gombméretek, és a túl hosszú űrlapok gyorsabban fárasztóvá válnak. Emiatt a reszponzív működés nem pusztán formai kérdés volt, hanem közvetlenül befolyásolta a használhatóságot.

Az \path{index.html} esetében mobilon különösen fontos, hogy a felhasználó ne vesszen el a különböző küldési lehetőségek között. Ha az üzenetküldés, sablonküldés, kérdőívindítás, médiaküldés és időzítés egyszerre túl nagy hangsúlyt kapna, a felület zsúfolttá válna. A használhatóság szempontjából ezért az volt a lényeges, hogy a kiválasztott művelethez kapcsolódó mezők kerüljenek előtérbe, és a felhasználó ne érezze úgy, hogy minden küldési módot egyszerre kell értelmeznie.

A \path{conv.html} oldalon a reszponzív működés még érzékenyebb kérdés. Egy beszélgetésnézet akkor használható jól, ha a tényleges üzenetfolyam kapja a legtöbb helyet. Kisebb kijelzőn a túl magas fejléc, a túl nagy navigációs elem vagy a sok helyet elfoglaló alsó sáv könnyen elveszi a teret attól, ami miatt a felhasználó megnyitotta az oldalt: a beszélgetés áttekintésétől. Ezért ebben a nézetben a reszponzív tervezés lényege az volt, hogy a chat-szerű élmény kisebb képernyőn se essen szét, az üzenetek pedig görgethető, olvasható formában maradjanak.

A beszélgetésnézet fejlesztése közben különösen jól látszott, hogy egy apró elrendezési döntés is erősen befolyásolja a használhatóságot. Amikor a fejléc vagy az üzenetküldő sáv túl nagy helyet foglalt el, a tényleges beszélgetési terület szűknek és nehezen használhatónak hatott. Emiatt a \path{conv.html} esetében törekedtem arra, hogy a fejléc és az alsó beviteli rész minél visszafogottabb legyen, a fő hangsúly pedig az üzenetfolyamon maradjon.

A \path{filter.html} esetében más probléma jelentkezik. A szűrési felület több mezőt is tartalmazhat, és ezek asztali nézetben még kényelmesen áttekinthetők. Mobilon viszont ugyanaz a mezőmennyiség gyorsan túl hosszúvá vagy nehezen kezelhetővé válhat. Ezért a szűrésnél fontosabbá vált a mezők logikus sorrendje és a fokozatos kitöltés lehetősége. A felhasználó először megadhatja a legfontosabb feltételeket, majd az eredményeket külön blokkban értelmezheti. Ez kevésbé látványos megoldás, mint egy mindent egyszerre mutató felület, de használat közben biztonságosabb.

A kérdőívszerkesztő oldalaknál a reszponzivitás azért volt különösen nehéz, mert a kérdések, válaszok és elágazások egymással kapcsolatban állnak. Asztali nézetben ezek az összefüggések könnyebben áttekinthetők, mobilon viszont gyorsan elveszhet a kapcsolat a kérdés és a hozzá tartozó válaszlogika között. Emiatt ezeknél az oldalaknál nem az volt a cél, hogy minden elem egyszerre látható maradjon, hanem az, hogy a szerkesztés egymás után követhető blokkokra bontva is érthető legyen. [R6], [R7], [R9]

\section{Validációs döntések és hibamegelőzés}

A validáció a WaccApp-ban nem önálló technikai extra, hanem a felhasználói élmény egyik alapja. A rendszer több olyan műveletet kezel, ahol egy hiányzó vagy hibás adat később sikertelen küldést, megszakadó kérdőívet vagy értelmezhetetlen szűrést okozhat. Emiatt a frontendnek már a művelet elindítása előtt jeleznie kell a nyilvánvaló problémákat.

Az \path{index.html} oldalon ez leginkább a küldési módok közötti különbségnél jelenik meg. Egy egyszerű szöveges üzenetnél a címzett és az üzenettartalom megléte a legfontosabb. Sablonos küldésnél már a sablon kiválasztása és a paraméterek kitöltése is szükséges. Médiafájl esetén a kiválasztott állomány megléte és típusa kerül előtérbe, időzített küldésnél pedig az is számít, hogy az időpont értelmezhető-e. A felületnek ezért nem általános hibát kell adnia, hanem az adott művelethez kapcsolódó problémát kell megmutatnia. [R14]

A kérdőívszerkesztő felületeken a validáció már nemcsak mezőszintű kérdés. Itt nem elég azt figyelni, hogy egy kérdés szövege vagy egy válaszmező ki van-e töltve. Azt is ellenőrizni kell, hogy a válaszokból kialakuló útvonal folytatható-e. Ha egy válaszhoz nem tartozik következő kérdés vagy lezárás, akkor a kérdőív futás közben megakadhat. Ez a felhasználó számára nem feltétlenül szerkesztési hibaként jelenik meg, hanem úgy, mintha a rendszer nem működne. Ezért fontos, hogy a hibás kapcsolatok lehetőleg már a szerkesztéskor láthatóvá váljanak.

A \path{filter.html} esetében a validáció más jellegű. Itt nem küldési hibát kell megelőzni, hanem azt, hogy a felhasználó értelmetlen vagy félrevezető keresést indítson. A dátumintervallum például csak akkor értelmezhető, ha a kezdő és záró dátum logikus sorrendben van. A név és telefonszám kapcsolata szintén érzékeny pont, mert egy rossz párosítás hibás következtetéshez vezethet. A szűrőfelület validációja ezért nemcsak technikai ellenőrzés, hanem az eredmények értelmezhetőségét is védi.

A hibamegelőzésnél fontos volt, hogy a felület ne büntesse a felhasználót. Egy hibaüzenet akkor jó, ha segít javítani, nem pedig csak megállítja a folyamatot. A WaccApp-ban ezért a validáció célja az volt, hogy a felhasználó minél hamarabb lássa, mit kell pótolnia vagy módosítania. Egy hiányzó telefonszám, üres üzenet, rosszul beállított időpont vagy félbehagyott kérdőívkapcsolat esetén jobb azonnal visszajelzést adni, mint megvárni, amíg a backend vagy a külső platform utasítja vissza a kérést. [R13], [R14]

\section{Állapotjelzések és visszacsatolási logika}

A WaccApp több művelete aszinkron módon működik, ezért a felhasználó nem mindig kap azonnali tartalmi eredményt. Egy üzenetküldés, fájlfeltöltés, szűrés vagy beszélgetésbetöltés közben a háttérben kérés indul a backend felé, majd a válasz alapján frissül a felület. Ha ilyenkor semmilyen jelzés nem jelenik meg, a felhasználó könnyen azt hiheti, hogy nem történt semmi, és újra megnyomja a gombot, vagy megszakítottnak gondolja a műveletet.

Az üzenetküldésnél ez különösen fontos. A küldés gomb megnyomása után a felhasználónak rövid időn belül érzékelnie kell, hogy a rendszer feldolgozza a kérést. Nem kell bonyolult animáció vagy túlmagyarázott üzenet, de valamilyen állapotjelzés szükséges. Egy státuszszöveg, módosított gombállapot vagy rövid visszajelzés már elég lehet ahhoz, hogy a felhasználó tudja, hogy a kattintása nem veszett el.

A szűrésnél és a beszélgetésbetöltésnél a visszacsatolás más szerepet kap. Itt nem az a fő kérdés, hogy egy üzenet sikeresen elment-e, hanem az, hogy az adatok frissültek-e. A \path{filter.html} esetében a felhasználónak tudnia kell, hogy a keresés lefutott, akkor is, ha nincs találat. Egy üres eredménylista önmagában félreérthető lehet, mert nem derül ki belőle, hogy nincs adat, vagy a lekérdezés nem működött. A \path{conv.html} esetében hasonló a helyzet: ha a beszélgetés betöltése közben a felület üresnek látszik, az könnyen rendszerhibának tűnhet.

A médiafájloknál a visszajelzés azért kap nagyobb szerepet, mert a fájlküldés általában lassabb és érzékenyebb folyamat, mint egy egyszerű szöveg elküldése. A felhasználó először kiválaszt egy fájlt, majd várja, hogy az valóban feldolgozásra és küldésre kerüljön. Ilyenkor a felületnek nemcsak a végeredményt kell jeleznie, hanem azt is, hogy a rendszer dolgozik. Ellenkező esetben könnyen előfordulhat, hogy a felhasználó újra kiválasztja a fájlt vagy ismételten elindítja a küldést.

Az időzített műveleteknél a visszacsatolás még összetettebb, mert a művelet rögzítése és tényleges végrehajtása időben elválik egymástól. A felhasználó először azt várja, hogy a rendszer elfogadta-e az ütemezést, később pedig azt szeretné látni, hogy az adott küldés várakozik, sikeresen lefutott vagy hibára futott. Ezek az állapotok nem pusztán technikai információk, hanem a felhasználói kontroll részei. A WaccApp-ban ezért az állapotjelzés nem dekorációként jelenik meg, hanem a kommunikációs folyamat értelmezését segíti. [R13], [R14]

\section{Vizuális következetesség és alap hozzáférhetőség}

A WaccApp több különálló HTML-nézetből áll, ezért a vizuális következetesség különösen fontos volt. Ha minden oldal más gombhasználatot, más hibaüzenet-megjelenítést és más űrlaplogikát követne, a felhasználó minden nézetnél újra tanulná a rendszert. Ez nemcsak kényelmetlen lenne, hanem hibákhoz is vezethetne, főleg olyan műveleteknél, ahol üzenetküldésről, fájlcsatolásról vagy kérdőívindításról van szó.

A következetesség ugyanakkor nem azt jelenti, hogy minden oldalnak ugyanúgy kell kinéznie. Az \path{index.html}, a \path{conv.html}, a \path{filter.html} és a kérdőívszerkesztő oldalak más feladatot látnak el, ezért természetes, hogy az elrendezésük is eltér. A fontos az volt, hogy az azonos típusú elemek hasonló jelentést hordozzanak. Egy elsődleges műveleti gomb mindenhol fő cselekvést indítson, a hibaüzenetek a javítandó adatra irányítsák a figyelmet, a státuszjelzések pedig a folyamat aktuális állapotát mutassák.

Az alapvető hozzáférhetőségi szempontok a WaccApp-ban főként a mindennapi használhatóságot támogatták. A jól olvasható szöveg, a megfelelő kontraszt, a látható fókuszállapot és a kényelmesen használható gombméret nem látványos fejlesztési eredmény, mégis sokat számít. Egy túl kicsi gomb mobilon könnyen félrekoppintáshoz vezethet, egy gyenge kontrasztú hibaüzenet pedig pont akkor nem segít, amikor a felhasználónak szüksége lenne rá.

A rendszer nem teljes körű akadálymentesítési projektként készült, de a legfontosabb alapelveket figyelembe kellett venni. A felületnek olvashatónak, kezelhetőnek és érthetőnek kell maradnia akkor is, ha a felhasználó kisebb kijelzőn dolgozik, gyorsan akar üzenetet küldeni, vagy egy hosszabb kérdőíves szerkesztési folyamatban navigál. Ezek az apró döntések együtt határozzák meg, hogy a rendszer mennyire használható a gyakorlatban. [R6], [R7]

\section{A UX-döntések szerepe a WaccApp használhatóságában}

A WaccApp felhasználói élményét nem egyetlen látványos megoldás határozza meg, hanem sok kisebb döntés együttese. A többnézetes felépítés segít elkülöníteni a különböző feladatokat, a reszponzív működés használhatóbbá teszi a rendszert eltérő kijelzőméreteken, a validáció csökkenti a hibás műveleteket, az állapotjelzések pedig láthatóvá teszik a háttérben futó folyamatokat. Ezek együtt adják azt az élményt, hogy a felhasználó nem technikai akadályokkal, hanem kezelhető munkafolyamatokkal találkozik.

A fejezetben bemutatott döntések a korábbi fejezet követelményeire épülnek. Itt a kérdés az, hogy ezek a funkciók hogyan váltak használható felületi megoldássá. A küldési felületnél ez a megfelelő mezősorrendben és visszajelzésben jelenik meg. A beszélgetésnézetnél az időrendi és vizuálisan elkülönített megjelenítésben. A szűrésnél a feltételek kezelhető szervezésében. A kérdőívszerkesztésnél pedig abban, hogy a kérdések és válaszok kapcsolatai követhetők maradjanak. A WaccApp használhatósága tehát azon múlik, hogy a felhasználó mennyire tudja magabiztosan végig vinni a fő műveleteket. Ha érti, mit kell kitöltenie, látja, ha valami hiányzik, kap visszajelzést a folyamat állapotáról, és a különböző nézetek között váltva sem veszti el a rendszer logikáját, akkor a frontend teljesíti a szerepét. [R9], [R13], [R14]
